% 3-concluzii.tex starts here
\chapter{Concluzii}\label{ch:concl}

Lucrarea de față a avut ca obiectiv proiectarea și implementarea unei aplicații web full-stack moderne, menită să demonstreze integrarea coerentă a tehnologiilor actuale din ecosistemul TypeScript într-o arhitectură unitară.
Proiectul \textit{Kookio} a fost conceput atât ca aplicație funcțională, cât și ca studiu de caz pentru analiza arhitecturilor web hibride utilizate în contexte reale.

Adoptarea arhitecturii \textbf{PAN} (Prisma--Analog--Nx) a permis realizarea unui flux end-to-end tipizat, de la interfața utilizatorului până la nivelul de persistență al datelor.
Această abordare reduce discrepanțele dintre frontend și backend, elimină duplicarea contractelor de date și facilitează dezvoltarea incrementală a funcționalităților.
În acest context, arhitectura PAN se conturează ca un model viabil pentru aplicații web full-stack moderne.

Utilizarea \textbf{Angular} împreună cu \textbf{Analog} a permis implementarea unei aplicații cu randare server-side și hidratare client-side.
Această combinație îmbunătățește performanța percepută și compatibilitatea cu motoarele de căutare.
Integrarea logicii de server în aceeași aplicație cu interfața utilizatorului simplifică arhitectura generală și reduce necesitatea unui backend separat, fără a compromite separarea responsabilităților.

Pentru accesul la date, combinația dintre \textbf{Prisma} și \textbf{CockroachDB} oferă un model de persistență robust, compatibil cu standardele SQL și pregătit pentru scenarii de scalare viitoare.
Definirea declarativă a schemei bazei de date, alături de generarea automată a clientului TypeScript, contribuie la creșterea siguranței, lizibilității și mentenabilității codului.

Procesul de testare a evidențiat importanța combinării mai multor niveluri de validare.
Testele unitare și de integrare realizate cu \textbf{Vitest} asigură corectitudinea logicii interne a aplicației.
În completare, testele end-to-end cu \textbf{Playwright} validează funcționarea sistemului din perspectiva utilizatorului final, inclusiv în contextul randării server-side.

Din punctul de vedere al obiectivelor inițiale, proiectul îndeplinește cerințele funcționale ale unui MVP (Minimum Viable Product).
Aplicația oferă funcționalități de bază pentru gestionarea ingredientelor și demonstrează fezabilitatea arhitecturii propuse.
Structura modulară a monorepo-ului și utilizarea instrumentului Nx permit extinderea ulterioară a aplicației fără modificări structurale majore.

În ansamblu, studiul realizat prin aplicația \textit{Kookio} arată că arhitectura PAN poate fi utilizată ca model coerent pentru dezvoltarea aplicațiilor web full-stack moderne.
Aceasta oferă atât valoare practică, cât și un cadru solid pentru analiză academică și dezvoltare ulterioară.

\section*{Direcții viitoare}\label{sec:roadmap}

Direcțiile viitoare de dezvoltare includ extinderea mecanismelor de autentificare și autorizare, optimizarea performanței prin introducerea unor strategii de caching și integrarea unor servicii inteligente pentru recomandarea rețetelor.

Aceste extensii pot fi realizate incremental, fără a afecta fundamentele arhitecturale prezentate în această lucrare.
% 3-concluzii.tex ends here